\section{Examples}
\label{sec:examples}

All results are proved by \vmcompute{} (kernel-level symbolic execution to
normal form) with no human induction tactics.  The SC drives both sides of each
equation to the same residual; \texttt{reflexivity} closes the goal.

\subsection{List monad laws}

The List monad has $\mathsf{return}\ x = [x]$ and
$m \bind f = \mathsf{foldr}\ (\lambda x\ \mathit{acc}.\ \mathsf{append}\ (f\ x)\ \mathit{acc})\ [\,]\ m$.

\begin{theorem}[List monad laws, \texttt{MonadLaws.v}]
\[
  m \bind \mathsf{return} = m
  \qquad
  \mathsf{return}\ x \bind f = f\ x
  \qquad
  (m \bind f) \bind g = m \bind (\lambda x.\ f\ x \bind g)
\]
\end{theorem}

\noindent
The left-identity law is proved by a single driving step (unfold \texttt{bind}
on $[\,]$; drive \texttt{append} on $[\,]$; \texttt{reflexivity}).  The
right-identity law requires a cyclic backlink: driving through the list spine
produces a recursive call that matches the root.  The associativity law is the
hardest: the SC must commute two nested folds, discovering that
$\mathsf{append}\ (\mathsf{append}\ a\ b)\ c = \mathsf{append}\ a\ (\mathsf{append}\ b\ c)$
holds at each step --- itself a cyclic induction, handled by the SC's
trace condition.

\paragraph{Functor and coherence laws.}
We additionally prove $\map\ \mathsf{id}\ l = l$, $\map\ f\ (\map\ g\ l) =
\map\ (f \circ g)\ l$, and the coherence law $\map\ f\ l = l \bind
(\mathsf{return} \circ f)$ --- six theorems total.

\subsection{Maybe monad laws}

The Maybe monad has $\mathsf{return}\ x = \mathsf{just}\ x$ and
$\mathsf{bind\_maybe}\ \mathsf{nothing}\ f = \mathsf{nothing}$,
$\mathsf{bind\_maybe}\ (\mathsf{just}\ x)\ f = f\ x$.

\begin{theorem}[Maybe monad laws, \texttt{IndexCalc.v}]
All three monad laws hold for the Maybe monad on \texttt{Nat}.
\end{theorem}

\noindent
These are even simpler than the List laws: Maybe has no recursion, so all
driving terminates in two steps.  They serve as a calibration point.

\paragraph{Head-bind fusion.}
A compound example: $\mathsf{bind\_maybe}\ (\mathsf{safe\_head}\ l)\ f$
fuses to a single case analysis on $l$ (theorem \texttt{head\_bind\_fusion}).

\subsection{Insertion sort}

\begin{theorem}[\texttt{SortExamples.v}]
The following hold for insertion sort on \texttt{List Nat}:
\begin{align*}
  \length\ (\sort\ l) &= \length\ l \\
  \length\ (\ins\ x\ l) &= \mathsf{succ}\ (\length\ l) \\
  \mathsf{member}\ x\ (\ins\ x\ l) &= \mathsf{true} \\
  \mathsf{sum}\ (\sort\ l) &= \mathsf{sum}\ l \\
  \length\ (\sort\ (\sort\ l)) &= \length\ l
\end{align*}
\end{theorem}

\noindent
The most surprising result is $\mathsf{sum}\ (\sort\ l) = \mathsf{sum}\ l$:
the SC discovers that insertion sort is a permutation (it preserves the sum)
by driving through the comparison and proving that the two branches of the
\texttt{leb} case produce structurally equivalent accumulations.
$\mathsf{member}\ x\ (\ins\ x\ l) = \mathsf{true}$ requires reasoning
about \texttt{leb} being reflexive.

\subsection{Vector index normalisation}

The SC normalises dependent type indices.  We prove:
\[
  \mathsf{Vec}\ A\ (\length\ (\map\ f\ l)) \equiv \mathsf{Vec}\ A\ (\length\ l)
\]
as an exact equality of index terms after supercompilation (theorem
\texttt{vec\_map\_index\_normalises}).  This is the foundational step for
dependent vector operations: it shows that the SC can be used as a typechecker
for programs involving non-trivial index computations.

\subsection{Plus associativity}

$\mathsf{plus}\ m\ (\mathsf{plus}\ n\ k) = \mathsf{plus}\ (\mathsf{plus}\ m\ n)\ k$
is proved by the SC (theorem \texttt{plus\_assoc\_killed}).  The SC drives
through the left argument of \texttt{plusL} and discovers the same recursive
structure on both sides.  Commutativity ($\mathsf{plus}\ m\ n = \mathsf{plus}\
n\ m$) is \emph{not} proved --- it requires a strengthened induction hypothesis,
discussed in Section~\ref{sec:limits}.

\subsection{The $\omega$-rule: reverse-append fusion}
\label{sec:reverse-append}

The most intricate example demonstrates the full four-layer pipeline.
The identity is:
\[
  \mathsf{reverse}\ (\mathsf{append}\ \mathit{xs}\ \mathit{ys})
  = \mathsf{append}\ (\mathsf{reverse}\ \mathit{ys})\ (\mathsf{reverse}\ \mathit{xs})
\]

The standard SC (all three layers) terminates on both sides but produces
\emph{different} residuals --- it cannot fuse them (theorem
\texttt{std\_sc\_gets\_stuck} in \texttt{LLMLemmaInAction.v}:
\texttt{r\_rev\_std $\neq$ r\_append\_rev\_std}).

The LLM lemma proposer, given the function definitions (in scrambled form),
proposes:
\[
  \mathsf{revAcc}\ (\mathsf{append}\ \mathit{xs}\ \mathit{ys})\ \mathit{acc}
  = \mathsf{revAcc}\ \mathit{ys}\ (\mathsf{revAcc}\ \mathit{xs}\ \mathit{acc})
\]

The sub-SC validates this lemma by induction on $\mathit{xs}$
(\texttt{lemma\_rev\_acc\_append\_distrib\_validated}, \vmcompute{}).

With the lemma in the driving step, the lemma-driven SC produces
\emph{identical} residuals for both sides (\texttt{with\_lemma\_they\_fuse}),
and the fused residual is provably different from the standard SC output
(\texttt{lemma\_driven\_is\_better}).

\subsection{Other $\omega$-rule examples}

The same pattern yields two additional results:

\begin{itemize}[leftmargin=*]
\item \textbf{Sort correctness}: $\mathsf{sorted}\ (\mathsf{sort}\ l) = \mathsf{true}$
  with lemma $\mathsf{sorted}\ (\ins\ x\ l) = \mathsf{sorted}\ l$.
\item \textbf{Succ distribution}: $\mathsf{sum}\ (\map\ \mathsf{succ}\ l)
  = \mathsf{sum}\ l + \length\ l$ with lemma
  $\mathsf{succ}\ m + n = \mathsf{succ}\ (m + n)$.
\end{itemize}

\subsection{Summary}

\begin{center}
\small
\begin{tabular}{lcc}
\toprule
\textbf{File} & \textbf{Theorems proved} & \textbf{Admitted} \\
\midrule
\texttt{HardExamples.v}          & 8  & 0 \\
\texttt{SortExamples.v}          & 9  & 0 \\
\texttt{MonadLaws.v}             & 6  & 0 \\
\texttt{IndexCalc.v}             & 10 & 0 \\
\texttt{CanonRoundtrip.v}        & 16 & 0 \\
\texttt{OmegaExamples.v}         & 12 & 0 \\
\texttt{LLMLemmaInAction.v}      & 4  & 0 \\
\midrule
\textbf{Total}                   & \textbf{65} & \textbf{0} \\
\bottomrule
\end{tabular}
\end{center}
